\chapter{Family--defect KO6 quiver embedding and spectral-sign gate}

Moment-map gate построил frozen--gauged--frozen quiver, но оставил открытым
его статус как finite geometry. Здесь предъявляется явная KO6 matrix model,
проверяются order-zero и first-order conditions и затем сравнивается её
ordinary spectral quartic с требуемым moment-map action. Результат
разделяется: representation embedding проходит, ordinary
\(\Tr D_F^4\) route закрывается по знаку mixed quartic.

\section{Finite algebra и bimodule labels}

Минимальная real algebraic skeleton имеет вид
\begin{equation}
 \mathcal A_{\rm q}
 =\mathbb R_0\oplus M_3(\mathbb R)_G\oplus\mathbb C_2.
 \label{eq:family-quiver-finite-algebra}
\end{equation}
Particle chain использует bimodule labels
\begin{equation}
 (0,0)^{\oplus3}
 \xrightarrow{\ X\ }
 (G,0)
 \xrightarrow{\ Y\ }
 (G,2).
 \label{eq:family-krajewski-label-chain}
\end{equation}
Первый node является scalar representation с multiplicity три. Его
commutant даёт global family frame. Middle и right modules имеют один и тот
же left \(M_3(\mathbb R)\)-label, поэтому несут один gauged connected
\(SO(3)_G\), а не два независимых gauge factors.

Первый edge меняет left label при фиксированном right label и поэтому может
быть arbitrary matrix
\begin{equation}
 X\in M_3(\mathbb R).
\end{equation}
Второй edge меняет right label при фиксированном left label. Exact
first-order condition заставляет его коммутировать со всем
\(M_3(\mathbb R)\):
\begin{equation}
 [Y,A]=0\quad\forall A\in M_3(\mathbb R)
 \qquad\Longrightarrow\qquad
 Y=\Phi I_3.
 \label{eq:first-order-schur-pairing}
\end{equation}
Machine kernel dimension равен единице. Random nonscalar \(Y\) даёт
first-order residual \(2.28\), тогда как scalar branch имеет нулевой
residual.

\section{KO6 completion}

Particle grading равен \((+,-,+)\). Добавим conjugate transposed-label
chain с противоположным grading и определим \(J_F\) обменом двух chains с
complex conjugation. Получается complex Hilbert space dimension
\begin{equation}
 \dim_{\mathbb C}\mathcal H_F=18.
\end{equation}
Для
\begin{equation}
 D_p=
 \begin{pmatrix}
 0&X^\dagger&0\\
 X&0&\bar\Phi I_3\\
 0&\Phi I_3&0
 \end{pmatrix},
 \qquad
 D_F=D_p\oplus\bar D_p,
 \label{eq:family-quiver-ko6-dirac}
\end{equation}
прямая проверка даёт
\begin{align}
 D_F&=D_F^\dagger,\\
 \{D_F,\Gamma_F\}&=0,\\
 [D_F,J_F]&=0,\\
 \{J_F,\Gamma_F\}&=0.
\end{align}
Все residuals равны нулю в double precision. Для полного basis algebra
также выполнены
\begin{align}
 [\pi(a),J\pi(b)J^{-1}]&=0,\\
 [[D_F,\pi(a)],J\pi(b)J^{-1}]&=0.
\end{align}

\begin{theorem}[KO6 quiver representation pass]
Chain~\eqref{eq:family-krajewski-label-chain} имеет explicit
18-dimensional self-adjoint odd KO6 completion, проходит order-zero и
first-order conditions и автоматически фиксирует pairing edge как
\(Y=\Phi I_3\). Поэтому moment-map quiver совместим с finite Krajewski
representation на algebraic level.
\end{theorem}

Classification logic сверена с T.~Krajewski,
\emph{Classification of Finite Spectral Triples},
\path{arXiv:hep-th/9701081}. Ограничения scalar potentials обычным spectral
action сверены с \path{arXiv:hep-th/9803199}; общий spectral principle ---
с \path{arXiv:hep-th/9606001}.

\section{Ordinary spectral quartic}

Теперь проверим, совпадает ли ordinary spectral action этого же
\(D_F\) с moment-map norm. На radial slice
\begin{equation}
 X=\rho I_3,
 \qquad
 \Phi=r\in\mathbb R
\end{equation}
particle half-traces равны
\begin{align}
 \Tr D_p^2&=6(\rho^2+r^2),\\
 \Tr D_p^4&=6(\rho^2+r^2)^2.
 \label{eq:family-quiver-ordinary-traces}
\end{align}
Следовательно, mixed quartic coefficient равен
\begin{equation}
 +12\rho^2r^2.
\end{equation}
Moment-map norm, напротив, имеет вид
\begin{equation}
 \tau_3(\mu_G^2)
 =(\rho^2-r^2)^2
 =\rho^4-2\rho^2r^2+r^4.
 \label{eq:family-quiver-moment-sign}
\end{equation}
Требуемый mixed coefficient отрицателен. Добавление quadratic term
\(-b\Tr D_p^2\) не меняет quartic sign, поэтому любой ordinary one-trace
polynomial \(a\Tr D_p^4-b\Tr D_p^2\) сохраняет эту обструкцию.

Попытка использовать grading также не помогает:
\begin{equation}
 \Str D_p^4=0.
\end{equation}
Зануление суперследа вызвано попарной симметрией нечётного спектра, а не
требуемой разностью матриц Грама для входящих и исходящих стрелок.

Все trace identities проверены на нескольких radial points; residuals
меньше \(7\times10^{-16}\).

\begin{nogocriterion}[Ordinary spectral mixed-sign no-go]
Для explicit KO6 quiver ordinary quartic \(\Tr D_F^4\) видит
\((\rho^2+r^2)^2\), тогда как required moment map видит
\((\rho^2-r^2)^2\). Ни quadratic spectral term, ни graded supertrace не
исправляют mixed sign. Поэтому moment-map condensate нельзя объявлять
следствием стандартного one-trace spectral action этого \(D_F\).
\end{nogocriterion}

\section{Что именно остаётся открытым}

No-go не уничтожает quiver action: в supersymmetric/quiver gauge theory
moment map является curvature auxiliary \(D\)-sector, а не polynomial
только от self-adjoint \(D_F=d+d^\dagger\). Project также уже располагает
relative-curvature precedent, где conditional expectation удаляет common
diagonal curvature и оставляет norm нужного oriented block.

Поэтому остаются три существенно различные routes:
\begin{enumerate}
 \item auxiliary \(D\)-field, чья algebraic elimination даёт
 \(\tau_3([d,d^\dagger]_G^2)\);
 \item relative/mapping-cone curvature action, канонически выделяющее
 middle-node incoming-minus-outgoing block;
 \item более общий non-ordinary spectral functional с заранее выведенной
 polarization, но без fitted counterterm.
\end{enumerate}

Первый и второй routes согласованы с quiver moment-map literature
\path{arXiv:hep-th/0405230}, \path{arXiv:0807.4734} и graded quiver gauge
construction \path{arXiv:0905.2338}. Но project должен вывести один из них
из собственной finite geometry и получить ту же kinetic normalization.

\begin{protocol}[Relative-moment continuation]
Следующий gate должен построить canonical conditional expectation или
auxiliary elimination на algebra
\(\mathbb R_0\oplus M_3(\mathbb R)_G\oplus\mathbb C_2\), доказать
\begin{equation}
 S_{\rm rel}=\tau_3([d,d^\dagger]_G^2),
\end{equation}
и проверить, что общий trace weight относительно \(|D X|^2\) и
\(|D\Phi|^2\) равен единице. Только после этого normalized condensed saddle
предыдущей главы станет spectral-parent result.
\end{protocol}

Таким образом, representation gap закрыт, но ordinary spectral-action route
закрыт отрицательно. Открыт более узкий relative/auxiliary curvature gate.

Следующий audit закрывает ordinary mapping-cone перенос и strict
\(SO(3)\)-adjoint D-term, но находит coefficient-free Legendre realization
на self-adjoint module \(\operatorname{Sym}_3(\mathbb R)\). Поэтому новый
разрыв состоит уже в выводе этого module из degree-two junk quotient или BV
complex той же KO6 geometry.

Повторная литературная проверка закрывает naive real auxiliary reading.
Canonical mapping-cone norm сохраняет только endpoint product; strict
\(SO(3)\) moment-map pairing исчезает на real slice; real positive
auxiliary minimum даёт quartic противоположного знака. Открыты только
positive degree-two junk-quotient curvature либо imaginary
Hubbard--Stratonovich quantum measure с заранее выведенным contour.

Машинный ledger сохранён в
\path{s2t_v4_family_defect_ko6_quiver_embedding_gate_results.json};
генератор ---
\path{s2t_v4_family_defect_ko6_quiver_embedding_gate.py}.
